\documentclass[11pt]{article}
% Shared preamble for all ML-RESEARCH papers.
% Papers load it with: % Shared preamble for all ML-RESEARCH papers.
% Papers load it with: % Shared preamble for all ML-RESEARCH papers.
% Papers load it with: \input{../../../papers-common/preamble.tex}
\usepackage[utf8]{inputenc}
\usepackage[T1]{fontenc}
\usepackage{lmodern}
\usepackage{amsmath,amssymb,amsthm}
\usepackage{graphicx}
\usepackage{booktabs}
\usepackage{hyperref}
\usepackage[margin=1.1in]{geometry}

\newtheorem{theorem}{Theorem}
\newtheorem{proposition}{Proposition}
\newtheorem{definition}{Definition}
\newtheorem{remark}{Remark}

\newcommand{\R}{\mathbb{R}}
\newcommand{\E}{\mathbb{E}}
\newcommand{\Prob}{\mathbb{P}}
\newcommand{\ip}[2]{\langle #1, #2 \rangle}

\usepackage[utf8]{inputenc}
\usepackage[T1]{fontenc}
\usepackage{lmodern}
\usepackage{amsmath,amssymb,amsthm}
\usepackage{graphicx}
\usepackage{booktabs}
\usepackage{hyperref}
\usepackage[margin=1.1in]{geometry}

\newtheorem{theorem}{Theorem}
\newtheorem{proposition}{Proposition}
\newtheorem{definition}{Definition}
\newtheorem{remark}{Remark}

\newcommand{\R}{\mathbb{R}}
\newcommand{\E}{\mathbb{E}}
\newcommand{\Prob}{\mathbb{P}}
\newcommand{\ip}[2]{\langle #1, #2 \rangle}

\usepackage[utf8]{inputenc}
\usepackage[T1]{fontenc}
\usepackage{lmodern}
\usepackage{amsmath,amssymb,amsthm}
\usepackage{graphicx}
\usepackage{booktabs}
\usepackage{hyperref}
\usepackage[margin=1.1in]{geometry}

\newtheorem{theorem}{Theorem}
\newtheorem{proposition}{Proposition}
\newtheorem{definition}{Definition}
\newtheorem{remark}{Remark}

\newcommand{\R}{\mathbb{R}}
\newcommand{\E}{\mathbb{E}}
\newcommand{\Prob}{\mathbb{P}}
\newcommand{\ip}[2]{\langle #1, #2 \rangle}


\title{The Maximum-Likelihood Principle}
\author{Yuval Lavie}
\date{\today}

\begin{document}
\maketitle

\begin{abstract}
Maximum-likelihood estimation answers one question: \emph{among a family of
candidate distributions, which one makes the observed data least surprising?}
We state the principle, derive the closed-form estimators for the Bernoulli,
Multinoulli, and Gaussian families, and note the properties that make MLE the
default estimator of statistics. The experimental section verifies, on
tracked runs, that gradient ascent on the log-likelihood
(\texttt{mlr.models.torch.mle}) recovers the same estimates as the analytic
formulas (\texttt{mlr.models.pure.mle}).
\end{abstract}

\section{The principle}

Let $S = (x_1, \dots, x_n)$ be i.i.d.\ draws from an unknown distribution,
and let $\{p_\theta : \theta \in \Theta\}$ be a parametric family of
candidate densities (or mass functions). The \textbf{likelihood} of a
parameter is the probability the family assigns to the sample as a function
of $\theta$:
\begin{equation}
  L_S(\theta) \;=\; \prod_{i=1}^{n} p_\theta(x_i).
\end{equation}

\begin{definition}[Maximum-likelihood estimator]
$\hat{\theta}_{\mathrm{MLE}} \in \arg\max_{\theta \in \Theta} L_S(\theta)$.
\end{definition}

Because $\log$ is strictly increasing, the same maximizer is found on the
\textbf{log-likelihood} $\ell_S(\theta) = \sum_i \log p_\theta(x_i)$, which
turns the product into a sum and is what one differentiates in practice.
Equivalently, MLE \emph{minimizes} the average negative log-likelihood
(NLL)
\begin{equation}
  \widehat{L}_S(\theta) \;=\; -\tfrac{1}{n} \sum_{i=1}^{n} \log p_\theta(x_i),
  \label{eq:nll}
\end{equation}
which is the loss reported by every estimator in the library.

\begin{remark}[Why NLL is the right loss]
Writing $\hat{q}$ for the empirical distribution of the sample,
$\widehat{L}_S(\theta) = H(\hat{q}) + \mathrm{KL}(\hat{q} \,\|\, p_\theta)$
up to reorganization: minimizing NLL is minimizing the Kullback–Leibler
divergence from the empirical data distribution to the model family.
Cross-entropy losses in machine learning are this identity in disguise.
\end{remark}

\section{Three closed forms}

\begin{proposition}[Bernoulli]
For $x_i \in \{0,1\}$ with $p_\theta(x) = \theta^x (1-\theta)^{1-x}$,
$\;\hat{\theta} = \bar{x} = \tfrac{1}{n}\sum_i x_i$.
\end{proposition}
\begin{proof}
$\ell_S(\theta) = \big(\sum_i x_i\big) \log\theta + \big(n - \sum_i x_i\big)
\log(1-\theta)$. Setting $\partial_\theta \ell_S = 0$ gives
$\sum_i x_i / \theta = (n - \sum_i x_i)/(1 - \theta)$, i.e.\
$\theta = \bar{x}$; $\ell_S$ is strictly concave, so this is the maximum.
\end{proof}

\begin{proposition}[Multinoulli]
For $x_i \in \{1, \dots, k\}$ with class probabilities
$\theta_1, \dots, \theta_k$ summing to one,
$\;\hat{\theta}_j = n_j / n$ where $n_j = |\{i : x_i = j\}|$.
\end{proposition}
\begin{proof}
Maximize $\sum_j n_j \log \theta_j$ subject to $\sum_j \theta_j = 1$ with a
Lagrange multiplier: $\partial_{\theta_j} \big[ \sum_j n_j \log\theta_j -
\lambda \sum_j \theta_j \big] = n_j/\theta_j - \lambda = 0$, so
$\theta_j = n_j / \lambda$; the constraint forces $\lambda = n$.
\end{proof}

\begin{proposition}[Gaussian]
For $x_i \in \R$ with $p_{\mu,\sigma}(x) =
\tfrac{1}{\sqrt{2\pi}\sigma} e^{-(x-\mu)^2 / 2\sigma^2}$,
\begin{equation*}
  \hat{\mu} = \bar{x},
  \qquad
  \hat{\sigma}^2 = \tfrac{1}{n} \sum_i (x_i - \bar{x})^2 .
\end{equation*}
\end{proposition}
\begin{proof}
$\ell_S(\mu, \sigma^2) = -\tfrac{n}{2}\log(2\pi\sigma^2) -
\tfrac{1}{2\sigma^2}\sum_i (x_i - \mu)^2$. Stationarity in $\mu$ gives
$\hat{\mu} = \bar{x}$; substituting and differentiating in $\sigma^2$ gives
$\hat{\sigma}^2$ as stated.
\end{proof}

\begin{remark}[The biased variance is not a bug]
$\E[\hat{\sigma}^2] = \tfrac{n-1}{n}\sigma^2$: the MLE of the variance is
biased. Bessel's $\tfrac{1}{n-1}$ correction is a \emph{different}
estimator; maximum likelihood genuinely produces $\tfrac{1}{n}$, and the
library implements exactly that.
\end{remark}

\section{Properties worth remembering}

\textbf{Invariance.} If $\hat{\theta}$ is the MLE of $\theta$, then
$g(\hat{\theta})$ is the MLE of $g(\theta)$ for any function $g$ — which is
why the torch estimators may search over unconstrained parameterizations
($p = \sigma(\vartheta)$, $\sigma = e^{s}$) and read off the constrained
estimate afterward.
\textbf{Consistency.} Under regularity conditions
$\hat{\theta}_n \to \theta^\ast$ as $n \to \infty$ when the family contains
the truth.
\textbf{Asymptotic efficiency.} $\sqrt{n}(\hat{\theta}_n - \theta^\ast)$ is
asymptotically normal with the inverse Fisher information as covariance — no
consistent estimator does better.

\section{Closed form versus gradient ascent}

In the library a distribution is one class (\texttt{mlr.distributions})
owning exactly the objects of this note: its density's NLL~\eqref{eq:nll},
the closed-form estimate \texttt{.mle()}, and an unconstrained
parameterization of the NLL for numerical search. A single model
(\texttt{model:\ mle}) then chooses the route: \texttt{method:\ mle}
evaluates the formula, while \texttt{method:\ gradient} feeds the NLL to an
optimizer — through hand-derived NumPy gradients
(\texttt{backend:\ pure}, verified against finite differences in the test
suite) or through autograd (\texttt{backend:\ torch}). When the
log-likelihood is concave in the chosen parameterization — true for all
three families here — every route must find the same optimum, so agreement
between the \texttt{-closed} and \texttt{-gradient} rows of
Table~\ref{tab:estimates} is a machine check of the derivations.

\begin{table}[h]
  \centering
  \small
  \input{../../results/tables/mle-estimates.tex}
  \caption{Latest tracked run per experiment
  (\texttt{research/mle/experiments/}). \texttt{-closed} rows evaluate the
  analytic formulas; \texttt{-gradient} rows run Adam on the NLL (torch
  autograd for Bernoulli/Multinoulli, pure NumPy gradients for the
  Gaussian). Held-out NLL and parameter estimates agree to the reported
  precision, and all recover the generating parameters.}
  \label{tab:estimates}
\end{table}

\begin{figure}[h]
  \centering
  \includegraphics[width=0.75\textwidth]{../../results/figures/gaussian-fit.pdf}
  \caption{Held-out view of the Gaussian experiment: sample histogram, the
  generating density, and the closed-form and gradient fits — the two fitted
  curves coincide.}
  \label{fig:gaussianfit}
\end{figure}

\end{document}
